\chapter{Introduction}
\label{ch:intro}

\section{Motivation}

Large language models can now articulate trading strategies with striking fluency. They
discuss momentum and mean-reversion, weigh news against technicals, and produce confident,
well-structured rationales for buying or selling. Yet a growing body of evidence shows that
this fluency is a poor guide to decision quality. LiveTradeBench~\cite{livetradebench2025}
finds that general LLM leaderboard scores barely correlate with realised trading returns ---
a model that ranks highly on generic benchmarks need not trade well --- which implies that
trading agents require their own, purpose-built evaluation. KellyBench~\cite{kellybench2026}
goes further, showing in a long-horizon sequential-decision setting that frontier models
articulate correct strategies but fail to execute them: they code a sizing rule and never
invoke it, they introduce label leakage into their own backtests, they declare success
prematurely, and their single-run results are unreliable. The recurring theme is a
\emph{knowledge--action gap}: stating the right thing and doing a different thing.

For an autonomous trading agent, the most dangerous instance of this gap is look-ahead bias.
Because the agent and its harness are built from code that has, in principle, access to the
entire historical record, it is alarmingly easy to leak future information into a past
decision --- and almost as easy to convince oneself, after the fact, that no such leak
occurred. The backtesting literature is unanimous that the only robust defence is
architectural: leakage must be impossible by construction, not avoided by discipline.

\section{Research question and contribution}

This thesis asks: \emph{can we build an autonomous trading agent on a standardised tool
protocol such that look-ahead is structurally impossible, and can we evaluate the quality of
its reasoning --- not merely its returns --- rigorously and by market regime?}

The contribution is twofold, and the order reflects where the academic value lies.

\begin{enumerate}
  \item \textbf{The \mcp-as-time-machine design.} The agent accesses market data, technical
  indicators, news sentiment, portfolio state, and order execution \emph{only} through Model
  Context Protocol~\cite{mcp2024} (\mcp) tool servers. A single simulation clock exposes
  \tnow, and every data tool refuses, at read time, to return any datum dated strictly after
  it. The agent calls the identical tools in backtest and in live operation; only the clock
  advances. Look-ahead is therefore impossible by construction (\cref{inv:nolookahead}), a
  guarantee pinned by tests written before the data layer and falsified by mutation testing.
  This is the central methodological contribution.
  \item \textbf{A regime-aware evaluation of reasoning quality.} On top of this substrate we
  measure three reasoning properties, each segmented by market regime: \emph{faithfulness}
  (does the executed trade match what a rational, policy-aware judge would do on the same
  evidence?), validated against blind human annotation; \emph{grounding} (are the numbers in
  the rationale real tool outputs?), computed deterministically; and a planned
  \emph{sophistication} rubric. Performance is reported against shifted, symmetrically-costed
  baselines with bootstrap confidence intervals.
\end{enumerate}

The primary, canonical result is a single trading agent. A multi-agent Portfolio Manager
--- analyst reports, a Bull/Bear research debate, a risk debate, and a synthesising allocator
in the manner of \textsc{TradingAgents}~\cite{tradingagents2024} --- is reported honestly as
a bonus extension. This hierarchy is deliberate: the project is graded on evaluation rigour,
not on the number of agents, and Agent Market Arena~\cite{ama2025} shows that agent
architecture, not backbone choice, drives behaviour --- so the interesting question about the
multi-agent system is behavioural (does deliberation close the single agent's
exit-discipline gap?), not whether more agents are intrinsically better.

\section{Why MCP, and why now}

The Model Context Protocol standardises how an LLM agent reaches external tools. Its
relevance here is not fashion but methodology: because the agent composes heterogeneous \mcp\
servers behind a uniform interface, and because each data server's wrapper owns the temporal
filter, the \emph{same} tool calls drive both backtest and live operation. This makes the
backtest/live equivalence exact rather than approximate, which is precisely what gives the
look-ahead guarantee its force. The \mcp\ layer is thus the architectural enabler of the
methodological claim, not a separable engineering detail.

\section{Scope}

The primary asset class is US equities at daily frequency over a window spanning at least two
market regimes, where regime labelling and baselines are cleanest and the closest comparable
benchmark, StockBench~\cite{stockbench2025}, is equity-centric. Cross-asset generalisation
(e.g.\ crypto) and reinforcement-learning fine-tuning of the agents are identified as future
work rather than attempted here, to keep the evaluation --- the graded core --- solid.

\section{Thesis structure}

\Cref{ch:related} situates the work against the evaluation benchmarks and multi-agent
frameworks it builds on. \Cref{ch:methodology} develops the methodological spine: the
\mcp-as-time-machine design, test-driven anti-leakage, the leakage channels a temporal filter
does not close, regime labelling, baselines, and the bootstrap. \Cref{ch:architecture}
describes the system: the \mcp\ servers, the single-agent orchestrator, the multi-agent
Portfolio Manager, the news corpus, and observability. \Cref{ch:evaluation} defines the
evaluation protocol and presents the rigour narrative of bugs caught and closed.
\Cref{ch:results} reports results. \Cref{ch:discussion} interprets them and states the
limitations honestly, and \cref{ch:conclusion} concludes.
